\section{Архитектура базы данных}
\headingtextsep

В разделе рассматривается структура данных системы, состав основных сущностей и связи между ними. Модель данных построена в нормализованном виде и ориентирована на соблюдение требований не ниже третьей нормальной формы.

\textheadingsep
\subsection{Ключевые сущности}
\headingtextsep

Ключевые сущности базы данных можно разделить на несколько групп: организационные сущности (`Workspace`, `WorkspaceMember`, `Project`), сущности управления работами (`TaskGroup`, `Task`, `TaskDependency`), сущности базы знаний (`Document`, `Asset`, `Note`), а также сущности агентного и интеграционного контуров (`Agent`, `AgentCredential`, `AgentSession`, `IntegrationConnection`, `AuditEvent`).

\textheadingsep
\subsection{Связи и ограничения}
\headingtextsep

В модели используются связи типов «один ко многим», «один к одному или многим» и «многие ко многим», последняя из которых реализуется через промежуточную сущность `AgentSessionTask`. Для обеспечения целостности применяются внешние ключи, а повторяющиеся и составные зависимости вынесены в отдельные таблицы, например `TaskDependency`, `AgentCredential` и `AuditEvent`.

\textheadingsep
\subsection{ERD-диаграмма}
\headingtextsep

ERD-диаграмма базы данных приведена в приложении А. На ней показаны основные сущности, их атрибуты, а также кардинальности и типы связей между таблицами.